% 34_body.tex — 산출물 ㉞ AI 판정 오류 사고 보고서 · Eval v2 · 오류 예산 개정 (빈 템플릿 / 완성 예시 공용 본문) — gen_product_templates.py 가 Product_specs.py 에서 생성
% 드라이버: 34_사고보고서Eval_v2_빈템플릿.tex (\wbblanktrue) / 34_사고보고서Eval_v2_완성예시.tex (\wbblankfalse — 워크북 §X.5 확정 후)
\documentclass[10pt,a4paper]{article}
\usepackage[a4paper,margin=16mm,top=14mm,bottom=20mm]{geometry}
\def\DocNum{34}\def\DocDay{7}\def\DocIdx{4}
\def\DocTitle{AI 판정 오류 사고 보고서 · Eval v2 · 오류 예산 개정}
\def\DocSub{Incident Report, Eval v2 \& Error Budget Revision}
\def\DocVariant{\B{빈 템플릿}{딥게이지 완성 예시}}
\usepackage{wbstyle}
\def\DocCirc{㉞}   % newunicodechar(㉛~㊵ 폴백) 뒤에 정의해야 활성 문자로 읽힌다
\begin{document}
\docheader
 {\Bg{[회사명 · 제품명]}{딥게이지 · GAUGE Inspect}}
 {\Bg{[INC-01 / D+520~540]}{INC-01 / 사고 D+520 · 보고서 D+525 · eval v2 D+540}}
 {\Bg{[작성자 — CTO·CS]}{오세진 CTO · 강지우 CS\rd{7mm}}}
 {\Bg{[승인자 — CEO]}{강민수\rd{7mm}}}

%% ------------------------------------------------------------------ 1
\secbar{1. 사고 개요}
\guide{일시·고객·사건(FN)·영향 로트 2,400개·청구 2,700만·해지 통보. 사실만, 판단은 §3.}
{\footnotesize
\begin{tabularx}{\linewidth}{|L{36mm}|Y|}\hline
\hd{항목} & \hd{내용} \\ \hline
\ifwbblank
\rd{9mm}사고 ID / 일시 &  \\ \hline
\rd{9mm}고객 / 라인 &  \\ \hline
\rd{9mm}무슨 일이 있었나(사실) &  \\ \hline
\rd{9mm}영향(수량·금액·고객 관계) &  \\ \hline
\rd{9mm}최초 인지 경로 &  \\ \hline
\else
사고 ID / 일시 & INC-01 / 2027-08-11(수) 15:40 \\ \hline
고객 / 라인 & [세영정공] 라인 2 — 교대 조, 검사원 1명(D+480 입사) \\ \hline
무슨 일이 있었나(사실) & 브래킷 외관 검사에서 초안 '양품'(신뢰도 0.87 ≥ 임계값 0.85). 조도 경고 표시. 검사원은 재촬영 안내를 건너뛰고(1탭) 승인. 17:10 로트 출하 \\ \hline
영향(수량·금액·고객 관계) & 로트 2,400개 H사 납품 → 크랙 발견·선별. 청구 2,700만(선별 1,100 + 운송 400 + 페널티 1,200) + 해지 통보(D+525). 계약 책임 한계 162만(ACV 1,620만 × 10\%) \\ \hline
최초 인지 경로 & H사 SQE → 세영 정미영 → 딥게이지(D+521 11:00 전화). 우리가 먼저 알지 못했다 — 라인 2 FN 신호 3개월 0건 \\ \hline
\fi
\end{tabularx}}

%% ------------------------------------------------------------------ 2
\secbar{2. 타임라인}
\guide{인지 → 내부 통보 → 고객 통보 → 임시 조치 → 원인 확인 → 재발 방지. 시각·담당·증빙.}
{\footnotesize
\begin{tabularx}{\linewidth}{|L{30mm}|Y|L{24mm}|L{40mm}|}\hline
\hd{시각} & \hd{일어난 일} & \hd{담당} & \hd{증빙} \\ \hline
\ifwbblank
\rd{8mm} &  &  &  \\ \hline
\rd{8mm} &  &  &  \\ \hline
\rd{8mm} &  &  &  \\ \hline
\rd{8mm} &  &  &  \\ \hline
\rd{8mm} &  &  &  \\ \hline
\rd{8mm} &  &  &  \\ \hline
\rd{8mm} &  &  &  \\ \hline
\rd{8mm} &  &  &  \\ \hline
\else
D+470 (06-15) & [세영] 박선재 '오후 3시 이후 크랙 놓친다' 구두 보고 & — & 증빙 없음 — 그것이 사실 \\ \hline
D+480 (06-25) & 라인 2 신입 검사원 입사 · 교육 1회 & 이수민 & 교육 기록 \\ \hline
08-11 15:40 & 초안 '양품' 0.87 · 조도 경고 · 재촬영 생략 · 승인 & 신입 검사원 & 판정 로그 B-12 \\ \hline
08-11 17:10 & 로트 2,400개 출하 & 세영 출하팀 & 출하 전표 \\ \hline
08-12 09:30 & H사 수입검사 크랙 발견 → 선별 → 통보 & [H사] 김도윤 & H사 통보서 \\ \hline
08-12 11:00 & 정미영 → 강민수 전화, 로그 요청 & 강민수 & 통화 기록 \\ \hline
08-12 16:00 & 임시 조치 — 오후 3시 이후 초안 OFF · 자동판정 ON 7사 OFF 권고 & 오세진 & 공지 메일 \\ \hline
08-16 / 08-31 & 청구 2,700 + 해지 통보 / eval v2 & 강민수 / 오세진 & 공문 / eval 리포트 v2 \\ \hline
\fi
\end{tabularx}}

%% ------------------------------------------------------------------ 3
\secbar{3. 근본 원인 분석 (8D / 5 Whys)}
\guide{기술 원인(조도·임계값·모델) / 절차 원인(검사원 회피 행동·로그 미감시) / 조직 원인(구두 보고가 기록되지 않음 — 박선재의 6월 발언). 원인 칸에 사람 이름을 쓰지 않는다 — 구조를 쓴다.}
{\footnotesize
\begin{tabularx}{\linewidth}{|L{8mm}|L{34mm}|Y|L{44mm}|}\hline
\hd{D} & \hd{단계} & \hd{내용} & \hd{근거} \\ \hline
\ifwbblank
\rd{8mm}D1 &  &  &  \\ \hline
\rd{8mm}팀 &  &  &  \\ \hline
\rd{8mm}D2 &  &  &  \\ \hline
\rd{8mm}문제 기술 &  &  &  \\ \hline
\rd{8mm}D3 &  &  &  \\ \hline
\rd{8mm}임시 조치 &  &  &  \\ \hline
\rd{8mm}D4 &  &  &  \\ \hline
\rd{8mm}근본 원인(5 Whys) &  &  &  \\ \hline
\rd{8mm}D5 &  &  &  \\ \hline
\rd{8mm}영구 조치 선정 &  &  &  \\ \hline
\rd{8mm}D6 &  &  &  \\ \hline
\rd{8mm}실행·검증 &  &  &  \\ \hline
\rd{8mm}D7 &  &  &  \\ \hline
\rd{8mm}재발 방지(시스템) &  &  &  \\ \hline
\rd{8mm}D8 &  &  &  \\ \hline
\rd{8mm}종결 &  &  &  \\ \hline
\else
D1 & 팀 & 오세진(리드) · 강지우 · 이수민 · 윤하경 · 외부: [세영] 정미영 & — \\ \hline
D2 & 문제 기술 & 오후 조도에서 크랙 부품을 초안 '양품'(0.87)으로 표시, 검사원이 재촬영을 건너뛰고 승인, 라인 FN 신호 3개월 0건이라 아무도 보지 않았다 & 로그 · ㉗ 항목 2 \\ \hline
D3 & 임시 조치 & 오후 초안 OFF · 자동판정 OFF 권고 · 세영 전 라인 초안 OFF & 08-12 공지 \\ \hline
D4 & 근본 원인 — 기술 & 오후 조도에서 신뢰도 과대(오전 학습 보정) → 0.85 통과. DS-v1이 오전만 대표(여집합 첫 줄) & X-02 · ㉗ 항목 2 \\ \hline
D4 & 근본 원인 — 절차 & 재촬영 생략 1탭·로그에 남지만 아무도 안 봄. FN 신호 0건 = '미감시'인데 리포트가 '여유'로 표시 & ㉗ 항목 5 근거 열 \\ \hline
D4 & 근본 원인 — 조직 & 구두 보고 기록 절차 없음(티켓은 팀장만). 신입 첫 4주에 초안 표시 & 정본 §2.3·§8.3 \\ \hline
D5 & 영구 조치 선정 & ① DS-v2 오후 세트 + 합격 전 오후 OFF ② 감시 대시보드(미감시·생략률) ③ 구두 보고 → CS 티켓 ④ 신입 첫 4주 초안 미표시 & 항목 8 \\ \hline
D6 & 실행·검증 & DS-v2 측정 / 첫 주간 리포트 / 티켓 3건 / 온보딩 체크리스트 & 항목 8 증빙 \\ \hline
D7 & 재발 방지(시스템) & 오류 예산 개정(항목 7) · HITL 재설계(생략 2탭 + 사유) · 여집합 조건 초안 표시 금지 & ㉗ v1.1 \\ \hline
D8 & 종결 & 세영 서면(항목 9) · 12월 갱신 시 결정 · 온담 제218문항의 답(Day8) & — \\ \hline
\fi
\end{tabularx}}

%% ------------------------------------------------------------------ 4
\secbar{4. 계약상 책임 vs 고객 주장}
\guide{책임 한계 162만(ACV 1,620만 × 10\%) vs 청구 2,700만. 고객 주장('시스템이 통과시킨 걸 믿었다')의 근거(GAUGE 로그)와 우리 HITL 문구. 협상 입장 3안.}
{\footnotesize
\begin{tabularx}{\linewidth}{|L{40mm}|Y|}\hline
\hd{항목} & \hd{내용} \\ \hline
\ifwbblank
\rd{10mm}계약 조항(책임 한계) &  \\ \hline
\rd{10mm}고객 주장과 근거 &  \\ \hline
\rd{10mm}우리 문구(㉗ 항목 6)가 방어하는 범위 &  \\ \hline
\rd{10mm}협상 입장 3안(전액/한계/중간)과 선택 &  \\ \hline
\rd{10mm}고객 관계 결정(유지·해지 수용) &  \\ \hline
\else
계약 조항(책임 한계) & '판정 초안의 오류로 인한 손해에 대한 당사 책임은 연 계약액의 10\%(162만 원)를 상한으로 한다.'(㉘ 항목 6 = ㉗ 항목 6) \\ \hline
고객 주장과 근거 & '우리 검사원은 시스템이 통과시킨 걸 믿었다' — 근거는 우리 로그(초안 양품 0.87·승인·조도 경고 플래그). 청구 2,700만 = G-36 \\ \hline
우리 문구(㉗ 항목 6)가 방어하는 범위 & 법적: 화면 문구 + 계약 상한 — 방어 가능. 관계적: 불가 — 경고를 띄우고도 초안을 표시한 설계, 무신호 3개월 방치. A-05b가 사고의 형태로 \\ \hline
협상 입장 3안(전액/한계/중간)과 선택 & ① 전액 — 관계·레퍼런스 / 선례 붕괴 ② 162만 — 선례 / 해지·채널 손실 ③ 중간(선택) — 162만 유지 + 6개월 무상 + 재발 방지 증빙 + 오후 OFF 공유. 사과와 책임 인정은 다른 문장 \\ \hline
고객 관계 결정(유지·해지 수용) & 해지 통보 → 12월 갱신까지 유예 요청, 유료 31에 포함(협상 중). 거절 시 수용, 레퍼런스 요청 안 함 \\ \hline
\fi
\end{tabularx}}

%% ------------------------------------------------------------------ 5
\secbar{5. Eval v2 결과와 트레이드오프}
\guide{v1 → v2: FN 8.1 → 2.4 / FP 14.6 → 21.3 / 단가 32 → 67원 / 지연 1.8 → 4.3초. FN을 낮춘 대가가 무엇인가 — 그리고 그것이 채택률에 미치는 인과를 적는다.}
{\footnotesize
\begin{tabularx}{\linewidth}{|L{30mm}|L{22mm}|L{22mm}|L{22mm}|Y|}\hline
\hd{지표} & \hd{v1} & \hd{v2} & \hd{변화} & \hd{이것이 무엇을 바꾸는가} \\ \hline
\ifwbblank
\rd{9mm}FN &  &  &  &  \\ \hline
\rd{9mm}FP &  &  &  &  \\ \hline
\rd{9mm}추론 단가 &  &  &  &  \\ \hline
\rd{9mm}지연 &  &  &  &  \\ \hline
\rd{9mm}채택률(㉮) &  &  &  &  \\ \hline
\else
FN & 24/296 = 8.1\% & 7/296 = 2.4\% & −5.7\%p(2-pass) & 불량 통과 1/3로. 그러나 오후 조도에서는 측정되지 않았다 — DS-v1 그대로 \\ \hline
FP & 173/1,184 = 14.6\% & 252/1,184 = 21.3\% & +6.7\%p & 양품 5장 중 1장이 불량으로. 재판정 부담 ↑. FP 예산(15\%) 초과 상태로 출발 \\ \hline
추론 단가 & 32원 & 67원 & ×2.1 & AI 원가 월 1,340만(순 SaaS MRR의 42.3\%). v1이면 약 640만. GM 49.8\%의 원인 절반 \\ \hline
지연 & 1.8초 & 4.3초 & ×2.4 & 라인 속도에서 4초는 '다음 부품' — 기다리지 않고 승인하거나 건너뛴다 \\ \hline
채택률(㉮) & (미측정) & 22\% — 24사 OFF & — & 인과: FN 2.4를 얻으려 FP 21.3을 감수 → 자동판정 ON 라인에서 양품 5장 중 1장이 재판정 요구 → OFF 요청 → 24사. 검사원 교육 문제가 아니라 FP의 결과. CTO는 FN을, CS는 채택률을 먼저 본다 — 같은 표 \\ \hline
\fi
\end{tabularx}}

%% ------------------------------------------------------------------ 6
\secbar{6. 합격 기준 개정 기록 — criteria drift 방지}
\guide{기준을 바꿨다면 ㉗ 항목 3의 절차를 통과했는가. 사고 후 개정의 정당화와 그 한계.}
{\footnotesize
\begin{tabularx}{\linewidth}{|L{36mm}|Y|}\hline
\hd{항목} & \hd{내용} \\ \hline
\ifwbblank
\rd{9mm}개정 전 → 후 &  \\ \hline
\rd{9mm}개정 근거 &  \\ \hline
\rd{9mm}승인권자·일자 &  \\ \hline
\rd{9mm}개정하지 않았어야 할 이유(반론) &  \\ \hline
\rd{9mm}데이터셋 변경 여부(모집단) &  \\ \hline
\else
개정 전 → 후 & FN ≤ 8.0\% \& OCR ≥ 90\% on DS-v1 → FN ≤ 3.0\% \& OCR ≥ 95\% on DS-v1(D+540) \\ \hline
개정 근거 & 사고 후 'FN 8\%는 월 n건 불량 통과 허용' — 2-pass가 2.4\%를 달성했으므로 상향 \\ \hline
승인권자·일자 & 오세진 CTO 전결, D+540 — ㉗ 항목 3 절차(높이는 개정 CTO 전결) 통과 \\ \hline
개정하지 않았어야 할 이유(반론) & ① 3.0 / 95는 2.4 / 95.2를 본 뒤의 숫자 — 결과에 맞춘 기준(drift) ② 데이터셋 그대로 — 사고 조건(오후 조도)은 어느 쪽에도 없다. 2.4\%는 오전에서의 값 ③ FP 21.3 > 예산 15인 모델을 합격 — 기준에 FP 없음. DS-v2가 먼저였어야 했다 \\ \hline
데이터셋 변경 여부(모집단) & 변경 없음 — 자동차 3세트 60건, 오전 9~11시. 여집합 첫 줄은 여전히 여집합 \\ \hline
\fi
\end{tabularx}}

%% ------------------------------------------------------------------ 7
\secbar{7. 오류 예산 정책 개정}
\guide{FP 예산 신설 여부 · 임계값 · 자동판정 기본값(ON/OFF) · 채택률 목표.}
{\footnotesize
\begin{tabularx}{\linewidth}{|L{36mm}|Y|Y|L{44mm}|}\hline
\hd{항목} & \hd{개정 전} & \hd{개정 후} & \hd{근거} \\ \hline
\ifwbblank
\rd{9mm}FN 예산 &  &  &  \\ \hline
\rd{9mm}FP 예산 &  &  &  \\ \hline
\rd{9mm}신뢰도 임계값 &  &  &  \\ \hline
\rd{9mm}자동판정 기본값 &  &  &  \\ \hline
\rd{9mm}소진 시 조치 &  &  &  \\ \hline
\else
FN 예산 & ≤ 8.0\%(라인별, HITL 수정 로그) & ≤ 3.0\% + FN 신호 0건 라인은 '미감시' 표기(무신호 ≠ 안전) & D4 절차 — 라인 2의 0건이 '여유'로 읽혔다 \\ \hline
FP 예산 & ≤ 15\% & ≤ 15\% 유지 — v2는 21.3\%로 초과 상태 출발. 초과 라인 자동판정 ON 불가 & 항목 5 채택률 인과 \\ \hline
신뢰도 임계값 & 0.85 & 0.90(오후 조도 사진 0.95) & 0.87이 통과한 사고 \\ \hline
자동판정 기본값 & 기능 없음 → D+410 기본 OFF(조건 미충족 출시) & 기본 OFF 고정 · ON은 FP ≤ 15\% 라인만 · B-17 재검토 조건 원복 & ㉖ Won't 조건 미충족 출시의 사후 정정 \\ \hline
소진 시 조치 & 라인 초안 OFF → 재학습 → 재측정 → ON & 유지 + 오후 조도 사진은 DS-v2 합격 전 OFF · 생략률 > 20\% 라인 팀장 알림 & D5 ① \\ \hline
\fi
\end{tabularx}}

%% ------------------------------------------------------------------ 8
\secbar{8. 재발 방지 조치 — 소유자·기한·증빙}
\guide{Day8 온담 218문항 중 제218문항이 이 표를 요구한다. 증빙이 없는 조치는 조치가 아니다.}
{\footnotesize
\begin{tabularx}{\linewidth}{|L{8mm}|Y|L{24mm}|L{22mm}|Y|L{18mm}|}\hline
\hd{\#} & \hd{조치} & \hd{소유자} & \hd{기한} & \hd{증빙(무엇으로 증명)} & \hd{상태} \\ \hline
\ifwbblank
\rd{9mm} &  &  &  &  &  \\ \hline
\rd{9mm} &  &  &  &  &  \\ \hline
\rd{9mm} &  &  &  &  &  \\ \hline
\rd{9mm} &  &  &  &  &  \\ \hline
\rd{9mm} &  &  &  &  &  \\ \hline
\rd{9mm} &  &  &  &  &  \\ \hline
\else
1 & DS-v2 오후 조도 세트 20건 수집·라벨·측정 & 오세진·라벨러 2 & D+620 & eval 리포트 v2.1(오후 세트 FN·FP 별도 행) & 진행 \\ \hline
2 & 감시 대시보드 — FN 신호·재촬영 생략률·'미감시' 표기 & 김태오 & D+610 & 첫 주간 리포트 발송 기록 & 진행 \\ \hline
3 & 오후 3시 이후 판정 초안 OFF(DS-v2 합격 전) & 오세진 & 완료(D+521) & 설정 로그 · 고객 공지 & 완료 \\ \hline
4 & 현장 구두 보고 → CS 티켓(검사원 직접 등록) & 강지우 & D+600 & 티켓 3건 이상 · 박선재 첫 티켓 & 완료(5건) \\ \hline
5 & 재촬영 생략을 2탭 + 사유로 · 생략률 주간 리포트 & 서하늘 & D+615 & 앱 릴리스 노트 · 생략률 리포트 & 진행 \\ \hline
6 & 신입 검사원 첫 4주 초안 미표시 · 온보딩 체크리스트 & 이수민 & D+600 & 체크리스트 v2 · 적용 로고 목록 & 완료 \\ \hline
\fi
\end{tabularx}}

%% ------------------------------------------------------------------ 9
\secbar{9. 고객 커뮤니케이션 문안}
\guide{세영정공에 보내는 서면 — 사실·책임 범위·조치·다음 단계. 사과와 책임 인정은 다른 문장이다.}
\begin{fieldbox}
\ifwbblank
\lines{7}{9.5mm}
\else
[세영정공] 정미영 품질팀장 귀하 (D+528)

(1) 사실 — 8월 11일 15시 40분 라인 2에서 당사 판정 초안이 크랙 부품을 '양품'으로 표시했고, 조도 경고가 표시된 상태에서 승인·출하되었음을 로그로 확인했습니다. 이런 일이 일어난 것을 깊이 유감으로 생각합니다.

(2) 책임 범위와 제안 — 당사의 책임은 계약 제n조(연 계약액의 10\%, 162만 원)를 상한으로 합니다. 별도로 향후 6개월 구독료를 청구하지 않고, 재발 방지 조치의 증빙을 월 1회 제출하겠습니다.

(3) 조치와 다음 단계 — 8월 12일부로 오후 3시 이후 판정 초안 표시를 중단했고, 오후 조도 데이터셋을 별도 구축해 합격 전까지 재개하지 않습니다. 라인별 감시 리포트를 매주 보내겠습니다. 박선재 검사원께서 6월에 말씀하신 내용이 기록되지 못한 점을 포함해, 현장 보고가 바로 등록되는 경로를 열었습니다.
\fi
\end{fieldbox}

\end{document}